\section{Démarche de transition écologique}

Ce projet s’inscrit pleinement dans une démarche de transition écologique en soutenant la connaissance et
la préservation de la biodiversité locale. Ce projet facilite la collecte d’images d’animaux sauvages pour
alimenter l’atlas de la biodiversité communale et aider les experts à mieux inventorier la faune du territoire.

% Ce projet s’inscrit pleinement dans une démarche de transition écologique en soutenant la connaissance, la valorisation et la préservation de la biodiversité locale. Il contribue activement à renforcer le lien entre les habitants, les acteurs locaux et leur environnement naturel, en favorisant une meilleure compréhension des écosystèmes présents sur le territoire.
% Grâce à la collecte d’images d’animaux sauvages, il facilite l’alimentation de l’atlas de la biodiversité communale et permet d’enrichir les bases de données existantes. Ces informations constituent une ressource précieuse pour les experts et les gestionnaires de la faune, leur permettant de mieux inventorier, suivre et protéger les espèces locales.
% Au-delà de l’aspect scientifique, le projet sensibilise également la population à l’importance de la faune et de la flore de proximité, incitant chacun à devenir acteur de la préservation du patrimoine naturel de sa commune.